\section{Metodología}

Para el desarrollo de la práctica, se llevo a cabo el siguiente procedimiento experimental:

\begin{enumerate}
    \item \textbf{Montaje y Verificación de Seguridad:}
    \begin{itemize}
        \item Verificar el correcto estado de los equipos e instrumentos de medición (VARIAC, vatímetro, osciloscopio, etc.).
        \item Realizar el montaje del circuito siguiendo el esquema de conexión, manteniendo el interruptor de prueba (breaker) en posición abierta hasta finalizar la revisión del cableado.
        \item Conectar el osciloscopio empleando resistencias shunt ($R_s$) para la visualización de las formas de onda, asegurando una correcta referencia a tierra.
    \end{itemize}

    \item \textbf{Ejecución del Ensayo con Carga Lineal:}
    \begin{itemize}
        \item Energizar el circuito ajustando el VARIAC hasta obtener la tensión nominal.
        \item Variar la carga resistiva en distintos niveles, registrando los valores de tensión ($V_2$), corriente ($A_2$) y potencia activa ($W_2$).
        \item Capturar las formas de onda de tensión y corriente mediante el osciloscopio para validar su comportamiento senoidal.
    \end{itemize}

    \item \textbf{Ejecución del Ensayo con Carga No Lineal:}
    \begin{itemize}
        \item Intercalar el puente rectificador en el lado secundario del transformador para alimentar la carga no lineal.
        \item Repetir la toma de mediciones de parámetros eléctricos, observando la distorsión inducida en la señal.
        \item Ajustar los elementos de la carga (inductor o arreglo capacitivo-resistivo) y documentar variaciones en el factor de potencia y el contenido armónico mediante instrumentos True RMS.
    \end{itemize}
\end{enumerate}